\chapter{Discussion and Threats to Validity}
\label{ch:discussion}

This chapter synthesizes the experimental findings, directly answers the four governing research questions, provides the theoretical foundation for why neural fact-checkers fail on structured assertions, and presents a rigorous four-part validity analysis.

% ******************************************************************************
\section{Synthesis: Answers to Research Questions}
\label{sec:rq_synthesis}
% ******************************************************************************

\subsection{RQ1: Efficacy of Modality-Based Routing}
\textbf{Question:} \textit{Does routing propositions into typed modalities outperform monolithic cross-encoders?}

\textbf{Answer:} Yes. HalDetect achieves \textbf{97.60\% selective accuracy} ($\text{MCC} = +0.9520$), outperforming the Non-Routed baseline ($60.00\%$, $\text{MCC} = +0.260$) by $+37.60\%$ on balanced ground truth ($p < 10^{-60}$). Modality routing enables domain-specific ALUs to resolve exact truth conditions without neural approximation.

\subsection{RQ2: Factual Precision in Structured Assertions}
\textbf{Question:} \textit{Can deterministic ALUs eliminate false alarms in quantitative and citation verification?}

\textbf{Answer:} Yes. Across 189 quantitative and 200 citation claims, the Dimensional SI ALU and CrossRef Registry Gateway achieve \textbf{99.47\%} and \textbf{100.00\% selective accuracy} respectively, with zero false alarms.

\subsection{RQ3: Generative Road-Test Generalization}
\textbf{Question:} \textit{How effectively does HalDetect catch real hallucinations across frontier foundation models?}

\textbf{Answer:} HalDetect achieves \textbf{100.0\% recall (31 / 31 errors caught)} across 1,798 generated propositions from seven foundation models, whereas Modern SOTA RAG and MiniCheck-7B achieve \textbf{0.0\% recall (0 / 31 caught)} due to lexical overlap bias.

\subsection{RQ4: Computational Footprint and Real-Time Feasibility}
\textbf{Question:} \textit{What are the compute and latency trade-offs compared to agentic RAG?}

\textbf{Answer:} HalDetect operates in \textbf{24.3 ms per claim on a consumer CPU} ($0\text{GB}$ GPU memory), running 24$\times$ faster than local GPU RAG ($578.5\text{ms}$) and 500$\times$ faster than cloud multi-search agents ($12,500\text{ms}$).

% ******************************************************************************
\section{Theoretical Analysis: Why Neural RAG Fails}
\label{sec:theoretical_failure}
% ******************************************************************************

The failure of neural cross-encoders to detect the 31 generative hallucinations stems from the \textbf{Attention Dilution Phenomenon}. When an LLM generates a numerical contradiction (e.g., \textit{"Apple M3 Max has 40 million transistors"} instead of $92\text{B}$), the continuous embedding distance between the claim and the retrieved reference article is dominated by matching context tokens (\textit{Apple, M3, Max, transistors, process}).

Because neural attention weights distribute over the high-frequency semantic context, the cross-encoder perceives overwhelming lexical overlap and assigns high entailment probability. In contrast, HalDetect isolates the numerical slot ($40\text{M} \neq 92\text{B}$) and executes discrete arithmetic discrepancy checks ($\Delta = 99.96\% > 5\%$).

% ******************************************************************************
\section{Threats to Validity}
\label{sec:threats_to_validity}
% ******************************************************************************

\subsection{Internal Validity}
Internal validity concerns whether confounding variables or data leakage influenced the measured performance. We addressed this through:
\begin{itemize}
    \item \textbf{AST Mechanical Leak-Check:} An automated Abstract Syntax Tree scanner (\texttt{\_leak\_check()}) verifies that predictor functions never access the \texttt{gold\_verdict} field during inference.
    \item \textbf{Decoupled Threshold Sweep:} Operating thresholds ($\tau_{\text{rel}} = 0.05, s > 0.60$) were tuned on an independent validation slice rather than fitted post-hoc to test sets.
\end{itemize}

\subsection{Construct Validity}
Construct validity examines whether our metrics genuinely reflect real-world hallucination detection. By reporting Matthews Correlation Coefficient (MCC) and Youden's J alongside selective accuracy, we ensure that performance is not inflated by class imbalance. Furthermore, evaluating probabilistic calibration via ECE ($8.88\%$) and Brier score ($0.0297$) confirms that confidence scores correspond to true error probabilities.

\subsection{External Validity}
External validity evaluates generalizability to unconstrained open-domain generation. Track 2 validated HalDetect across seven distinct foundation models of varying scales (7B to 120B) across six diverse domains, proving that the 5-tuple schema atomizer generalizes to varied prompt contexts.

\subsection{Limitations Inherited from Knowledge Bases}
HalDetect's coverage is bounded by the completeness of external knowledge graphs (Wikidata, World Bank, CrossRef). If an entity or macroeconomic indicator is absent from the structured graph, HalDetect abstains with an epistemic \texttt{UNVERIFIABLE} verdict.

% ******************************************************************************
\section{What the Evidence Does NOT Establish}
\label{sec:honest_bounds}
% ******************************************************************************

To maintain strict academic integrity, we explicitly state the boundaries of our findings:
\begin{enumerate}
    \item The evidence does \textbf{not} establish that neural cross-encoders are useless; rather, it proves they are poorly suited for arithmetic, interval containment, and citation verification.
    \item HalDetect does \textbf{not} eliminate the need for retrieval; it replaces unguided neural text matching with structured, typed knowledge retrieval.
    \item The 100\% hallucination recall observed in Track 2 applies to discrete factual confabulations (numbers, dates, DOIs) and does not guarantee complete detection of subtle rhetorical or aesthetic misalignments.
\end{enumerate}
